\chapter{Conclusions and Future Work}
\section{Conclusions}
This thesis has presented the development of a Digital Twin for a hand exoskeleton, following a complete and structured design pipeline that spans from the mechanical modeling of the device to the implementation and validation of a functional safety architecture. The work has been carried out within the ROS 2 framework, with the objective of providing a coherent, reproducible, and extensible platform for the design, simulation, and pre-validation of control and safety strategies prior to deployment on real hardware.

The main contributions of this thesis can be summarized as follows.

\textbf{Mechanical design and URDF modeling}. Starting from the reference dimensions provided in Bartalucci's thesis, the exoskeleton components were modeled parametrically in Fusion 360 and validated against the original design. The CAD assembly was converted into a URDF description using the ACDC4Robot plugin, and the resulting model was refined to handle the closed kinematic chains characteristic of the finger mechanism. The chain-opening strategy adopted — in which selected revolute joints are replaced by auxiliary reference frames whose coincidence is enforced at the software level — enabled the use of standard URDF-based tools while preserving the physical consistency of the original closed-loop design.

\textbf{Kinematic and dynamic modeling}. A constrained forward kinematics formulation was developed to enforce loop-closure constraints at runtime, producing kinematically admissible configurations for any value of the actuated crank angle. Building on this kinematic foundation, a reduced-order dynamic model was derived by projecting the full multibody dynamics onto the single actuated degree of freedom. The reduction exploits the constraint Jacobian to compute an admissible motion direction, onto which the inertia matrix, nonlinear terms, and all additional torque contributions — friction, passive biomechanical effects, external interaction forces — are projected. The resulting scalar nonlinear equation preserves the essential physics of the original system while enabling deterministic execution at 200 Hz within the ROS 2 environment. The open-loop validation confirmed the numerical stability of the kinematic closure solver, the correctness of the torque breakdown, and the physical consistency of the passive finger model.

\textbf{Control architecture}. A cascade admittance–trajectory control architecture was implemented and validated within the Digital Twin. The outer admittance loop generates a compliant motion reference in response to externally applied forces, encoding the desired interaction behavior through configurable virtual inertia, damping, and stiffness parameters. The inner trajectory loop tracks this reference using a PD controller with comprehensive model-based feed-forward compensation, achieving sub-degree tracking accuracy. The closed-loop validation through step and sinusoidal response experiments confirmed that the cascade architecture produces stable, overshoot-free transients consistent with the analytical predictions of the virtual dynamics.

\textbf{Functional safety architecture}. A dedicated functional safety layer was developed, organized into three cooperating layers: the ExoBridge enforcement node, which mediates all control signals and physically applies torque limits, velocity bounds, and stop conditions; the FaultMonitorMode detection plugin, which evaluates the system signals against configurable thresholds through a debounced, multi-level escalation logic; and the StateMachine governance node, which manages the safety state through a plugin-based architecture with fault latching, bridge supervision, and automatic downgrade capabilities. The architecture follows a defense-in-depth philosophy, in which redundant watchdogs at the bridge, monitor, and state machine levels ensure that no single point of failure can leave the system unprotected. A fault injection framework was developed in parallel to enable systematic validation of the safety mechanisms under controlled and reproducible fault scenarios.

\textbf{Model-based fault detection}. An observer module comprising four complementary detection layers — a Luenberger state observer, an admittance inversion estimator, a generalized momentum observer, and a rate-of-change guard — was implemented and integrated into the Digital Twin launch system. The four layers provide coverage across all critical signal channels with detection latencies ranging from 5 ms to approximately 1.5 seconds. The module is fully operational and has been qualitatively validated with the fault injection framework, although it has not yet been connected to the safety decision chain.

\textbf{Hardware prototype documentation}. The hardware platform of the hand exoskeleton was comprehensively documented, covering the electromechanical components, wiring architecture, motor driver configuration, and board-level software. The modifications introduced to the OEM software — including a robust auto-reconnection mechanism, thread-safe GUI updates, and offline velocity estimation methods — were motivated by the requirements of the Digital Twin integration and provide the interface specification for a future hardware-in-the-loop deployment.

Taken together, these contributions demonstrate that the Digital Twin paradigm can effectively support the full development cycle of a hand exoskeleton for rehabilitation: from mechanical design and physics-based simulation to control synthesis, safety verification, and hardware interface definition. The Digital Twin acts as a virtual testbench that enables early-stage evaluation of design choices, control parameters, and safety strategies in a safe and controlled environment, reducing the need for iterative hardware prototyping and minimizing the risk associated with real-world testing.

\section{Future Work}
Several directions for future development emerge naturally from the work presented in this thesis.
Hardware-in-the-loop integration. The most immediate next step is the integration of the ROS 2 Digital Twin stack with the physical prototype described in Chapter 7. The WebSocket-based distributed control architecture already developed by the industrial partner provides the communication layer between the Raspberry Pi on the exoskeleton and the ROS 2 workstation. The interface requirements — signal channels, sensor types, communication latencies, and actuation dynamics — have been fully documented in this thesis and constitute the specification against which the integration must be validated. This step would enable the first experimental comparison between the Digital Twin predictions and the actual hardware behavior, providing the data necessary for model parameter identification and calibration.

\textbf{Observer integration into the safety chain}. The model-based observer module described in Chapter 6 currently publishes diagnostic residuals but does not trigger any protective action. Integrating the observer into the functional safety state machine — through the SENSOR\_DEGRADED\_MODE state already reserved in the architecture — requires three well-defined tasks: systematic calibration of detection thresholds across the operating envelope, design of a decision logic (e.g., CUSUM or GLR-based) to convert continuous residuals into binary fault decisions, and validation of the integrated system under the full range of fault scenarios supported by the injection framework.

\textbf{Multi-finger extension}. The current Digital Twin models a single finger module. Extending the architecture to the full three-finger configuration of the prototype would require the instantiation of three parallel kinematic and dynamic chains sharing the same base frame, as well as the extension of the safety state machine to manage per-finger fault states. The parametric CAD design and the modular software architecture developed in this thesis are intended to facilitate this extension.

\textbf{Rehabilitation protocol implementation}. The admittance controller validated in this thesis implements a generic compliant interaction behavior. For clinical deployment, specific rehabilitation protocols — such as passive range-of-motion exercises, active-assisted movements, or resistance training — would need to be implemented as configurable exercise profiles within the admittance framework. The virtual stiffness and damping parameters could be programmed to vary over time according to a prescribed therapeutic trajectory, enabling adaptive rehabilitation sessions.

\textbf{Experimental validation of the safety architecture}. The functional safety mechanisms have been validated within the Digital Twin simulation environment using the fault injection framework. A comprehensive validation on the physical prototype — including the verification of the emergency stop chain, the bridge watchdog response times, and the actuator behavior under fault conditions — is essential before the system can be considered for clinical use. This validation should be conducted in accordance with the relevant safety standards for medical devices and rehabilitation robotics.